% === Chapter 2: Reward Modeling & Preference Optimization ===
\section{Thread 2: Reward Modeling \& Preference Optimization}
\label{sec:t2}

\begin{quote}
\textit{Core question: How can we transform vague, hard-to-formalize human preferences into training signals that can be optimized by gradient descent?}
\end{quote}

\noindent
One of the central challenges of post-training is that human judgments of ``good responses'' are holistic and comparative in nature,
rather than something that can be directly described by a simple loss function.
Humans can easily judge that ``response A is better than response B,'' but find it difficult to provide a precise scoring function.
This chapter traces the complete evolutionary path of this problem from formulation to (partial) resolution:
from the foundational RLHF framework (\S\ref{sec:t2-foundation}),
to the multiple branching research paths spawned by its core limitations (\S\ref{sec:t2-branches}),
to the recent convergence trends seeking new balance points between efficiency and effectiveness (\S\ref{sec:t2-convergence}),
and finally a discussion of open problems that remain unresolved (\S\ref{sec:t2-open}).

% ============================================================
\subsection{Problem Origin: Why Is SFT Not Enough?}
\label{sec:t2-origin}
% ============================================================

Supervised Fine-Tuning (SFT) trains models via maximum likelihood estimation (MLE) to imitate
human-annotated (instruction, response) pairs. While this paradigm is simple and effective (\crossref{1}{sec:t1}),
it suffers from two fundamental limitations.

\paragraph{Limitation 1: The mode-covering property of MLE and the absence of quality signals.}
The MLE objective $\mathcal{L}_{\text{SFT}} = -\mathbb{E}_{(x,y)\sim\mathcal{D}}\left[\log \pi_\theta(y|x)\right]$
treats all responses in the training set equally---regardless of quality, the model attempts to maximize their log-likelihood.
This leads to two consequences:
(1) the model learns the ``average level'' of the training data distribution rather than aligning with the best responses;
(2) when annotators vary in skill level, low-quality annotations directly degrade model performance,
and MLE has no intrinsic mechanism to distinguish and filter out such noise.

\paragraph{Limitation 2: The asymmetry of human judgment (judging $\neq$ generating).}
A deeper observation is that
human capabilities in ``generating high-quality responses'' and ``judging which response is better'' are asymmetric.
Generating an ideal response that satisfies all constraints (accurate, helpful, safe, fluent) is extremely difficult,
whereas comparing two existing responses is relatively easy and yields higher consistency.
Stiennon et al.\ \cite{stiennon2020summarize} found that inter-annotator agreement
on preference comparison tasks was approximately 77\%, significantly higher than on direct scoring tasks.
InstructGPT \cite{ouyang2022instructgpt} reported similar findings---
pairwise agreement among annotators was approximately 73\%.

\evolution{MLE can only imitate the ``average level''}{Leverage human comparison signals to distinguish quality}{Humans are better at ``judging'' than ``generating''}

\noindent
This observation forms the starting point for the entire field of preference optimization:
since humans are better at comparative judgment, we should design a training paradigm
that uses preference comparison signals (rather than direct demonstration signals) as the optimization objective.
Christiano et al.\ \cite{christiano2017rlhf} first proposed this idea in the reinforcement learning domain.
Ziegler et al.\ \cite{ziegler2020finetuning} were the first to transfer it to NLP---
demonstrating the feasibility of reward learning from comparisons on text continuation, summarization, and style transfer tasks,
and establishing a complete pipeline from human preferences to language model fine-tuning.
This paradigm was subsequently developed into a mature technical framework.

% ============================================================
\subsection{Foundational Method: The Classic RLHF Framework}
\label{sec:t2-foundation}
% ============================================================

The core idea of RLHF (Reinforcement Learning from Human Feedback) can be decomposed into two steps:
(1) training a reward model to learn human quality judgment criteria from preference comparisons;
(2) using this reward model as the optimization objective, adjusting the language model's policy through an RL algorithm (typically PPO).
This section traces three landmark works that brought this framework from inception to maturity.

% --- 2.2.1 ---
\subsubsection{From Games to Language: Reward Learning from Comparisons}
\label{sec:t2-christiano}

\landmark{christiano2017rlhf} first demonstrated on Atari games and MuJoCo continuous control tasks that
agents can learn complex behaviors solely from human preference comparisons over trajectory segments,
without access to the environment's true reward function.

\paragraph{Motivation and problem formalization.}
Traditional RL assumes the existence of an explicit reward function $r: \mathcal{S} \times \mathcal{A} \to \mathbb{R}$,
but in many practical scenarios (such as training a robot for dexterous manipulation, or training an agent to exhibit ``helpful'' behavior),
designing such a reward function is either extremely difficult or leads to reward hacking.
The core insight of Christiano et al.\ is that
although humans do not know the exact reward function,
they can judge which of two observed agent behaviors is better.

Formally, given two trajectory segments $\sigma^1 = (s_0^1, a_0^1, \ldots, s_k^1)$ and
$\sigma^2 = (s_0^2, a_0^2, \ldots, s_k^2)$,
human annotators select a preference label $\mu \in \{(1,0),\; (0,1),\; (0.5, 0.5)\}$.

\paragraph{Reward model training.}
Assuming human preferences follow the Bradley-Terry (BT) model \cite{bradley1952bt},
where preference probability is proportional to the exponential of the sum of rewards:
\begin{equation}
\label{eq:pref:bt-christiano}
P[\sigma^1 \succ \sigma^2] = \frac{\exp\!\sum_t \hat{r}(s_t^1, a_t^1)}{\exp\!\sum_t \hat{r}(s_t^1, a_t^1) + \exp\!\sum_t \hat{r}(s_t^2, a_t^2)}
\end{equation}
where $\hat{r}_\psi$ is a reward model parameterized by $\psi$ (a deep neural network).
The training objective is to minimize the cross-entropy loss:
\begin{equation}
\label{eq:pref:rm-loss-christiano}
\mathcal{L}_{\text{RM}} = -\mathbb{E}_{(\sigma^1, \sigma^2, \mu)}\left[
  \mu(1)\log P[\sigma^1 \succ \sigma^2] + \mu(2)\log P[\sigma^2 \succ \sigma^1]
\right]
\end{equation}

\paragraph{Policy optimization.}
After the reward model is trained, PPO \cite{schulman2017ppo} is used to optimize the policy $\pi_\theta$,
maximizing the reward signal provided by $\hat{r}_\psi$.
The entire system runs three processes asynchronously:
(1) the policy $\pi_\theta$ interacts with the environment to generate trajectories;
(2) human annotators provide preference judgments on trajectory segment pairs;
(3) the reward model $\hat{r}_\psi$ is continuously updated on preference data.

\paragraph{Key experimental results.}
In Atari tasks, using only approximately 5,500 human comparisons (about 1 hour of annotation),
the RLHF agent achieved performance comparable to reward shaping on some games;
in MuJoCo tasks (such as Hopper and Walker),
the RLHF agent successfully learned locomotion behaviors nearly identical to those of agents trained with the true reward.
These results established a key proposition:
\textbf{Human preference signals are sufficient to replace explicit reward functions for training RL agents.}

\begin{solutionbox}
\textbf{Core contribution}:
Established the complete framework for \textit{comparison-based reward learning}---
from the Bradley-Terry preference model to asynchronous reward model training to PPO policy optimization.
This framework became the prototype for all subsequent RLHF work.
\end{solutionbox}

% --- 2.2.2 ---
\subsubsection{From Games to Language: Learning to Summarize}
\label{sec:t2-summarize}

\landmark{stiennon2020summarize} applied the comparison-based reward learning framework of Christiano et al.\
for the first time in its entirety to a natural language generation task---text summarization---
and established the \textbf{SFT $\rightarrow$ RM $\rightarrow$ RL} three-stage pipeline that was subsequently widely adopted.

\paragraph{Three-stage pipeline.}

\textit{Stage 1: Supervised Fine-Tuning.}
GPT-3 series models (1.3B and 6.7B) were fine-tuned via SFT on the TL;DR dataset,
providing the models with basic summarization capability.

\textit{Stage 2: Reward Model Training.}
Human annotators provided preference judgments on pairs of model-generated summaries.
Given a post $x$ and two summaries $y_0, y_1$,
the reward model $r_\psi$ is trained with the following objective:
\begin{equation}
\label{eq:pref:rm-loss-summarize}
\mathcal{L}_{\text{RM}}(\psi) = -\mathbb{E}_{(x, y_w, y_l) \sim \mathcal{D}} \left[
  \log \sigma\!\left(r_\psi(x, y_w) - r_\psi(x, y_l)\right)
\right]
\end{equation}
where $y_w$ and $y_l$ denote the preferred (winning) and dispreferred (losing) responses respectively,
and $\sigma$ is the sigmoid function.
Note that unlike \cref{eq:pref:bt-christiano}, where rewards are summed over trajectory timesteps,
here the reward model directly scores complete $(x, y)$ pairs---
a key simplification that accompanied the migration of reward modeling from the RL domain to NLP.

\textit{Stage 3: RL Optimization.}
PPO is used to optimize the policy $\pi_\theta$, with the following objective:
\begin{equation}
\label{eq:pref:rl-obj-summarize}
\max_{\pi_\theta}\; \mathbb{E}_{x \sim \mathcal{D},\; y \sim \pi_\theta(\cdot|x)} \left[
  r_\psi(x, y) - \beta \log \frac{\pi_\theta(y|x)}{\pi_{\text{ref}}(y|x)}
\right]
\end{equation}
where $\pi_{\text{ref}}$ is the model obtained from the SFT stage (serving as the reference policy),
and $\beta$ is the KL penalty coefficient, used to prevent $\pi_\theta$ from diverging too far from $\pi_{\text{ref}}$---
a design that directly preempted the reward overoptimization problem that was later extensively studied (\S\ref{sec:t2-reward-hacking}).

\paragraph{Key findings.}
(1) A 6.7B model trained with human preference signals produced summaries of higher quality than a larger model trained with SFT alone,
demonstrating that \textbf{preference signals can compensate for gaps in model scale}.
(2) The reward model achieved an agreement rate of approximately 77\%, and reward model scores were highly correlated with human judgments,
validating the effectiveness of the Bradley-Terry model as a tool for modeling human preferences.
(3) A clear trade-off relationship exists between RM score and KL divergence:
as the KL penalty decreases, RM scores continue to rise, but the quality as assessed by actual human evaluation begins to decline past a certain point---
an early observation of the reward overoptimization phenomenon.

% --- 2.2.3 ---
\subsubsection{InstructGPT: Industrial-Scale Deployment of RLHF}
\label{sec:t2-instructgpt}

\landmark{ouyang2022instructgpt} scaled the three-stage pipeline described above to the general instruction-following setting,
providing the first production-scale demonstration of RLHF's effectiveness,
and establishing the standard paradigm for virtually all subsequent LLM post-training.

\paragraph{Industrial improvements to the pipeline.}

\textit{SFT stage}:
GPT-3 (175B) was trained on human-annotated demonstration data,
where annotators were asked to write ``ideal responses'' for user prompts.

\textit{RM stage}:
For each prompt, $K$ model outputs were sampled ($K \in \{4, 9\}$),
and annotators ranked all $\binom{K}{2}$ pairs.
The reward model training objective is:
\begin{equation}
\label{eq:pref:rm-loss-instructgpt}
\mathcal{L}_{\text{RM}}(\theta) = -\frac{1}{\binom{K}{2}} \mathbb{E}_{(x, y_w, y_l) \sim \mathcal{D}} \left[
  \log \sigma\!\left(r_\theta(x, y_w) - r_\theta(x, y_l)\right)
\right]
\end{equation}
Compared with \cref{eq:pref:rm-loss-summarize}, the key improvement is that
\textbf{all comparison pairs under the same prompt share a single forward pass},
substantially reducing the computational cost of training.
The RM was initialized from a 6B GPT-3 model (rather than 175B),
a design choice reflecting an important engineering principle:
\textbf{the reward model need not be the same scale as the policy model}.

\textit{PPO stage}:
The optimization objective adds pretraining loss mixing on top of \cref{eq:pref:rl-obj-summarize}:
\begin{equation}
\label{eq:pref:ppo-ptx}
\max_{\pi_\theta}\; \mathbb{E}_{x \sim \mathcal{D}_{\text{pref}},\; y \sim \pi_\theta}
  \left[r_\psi(x, y) - \beta \,\text{KL}\!\left[\pi_\theta \| \pi_{\text{ref}}\right]\right]
  + \gamma\,\mathbb{E}_{x \sim \mathcal{D}_{\text{pretrain}}}\left[\log \pi_\theta(x)\right]
\end{equation}
where the second term (PPO-ptx) mixes a portion of pretraining data into PPO training,
mitigating the alignment tax---the degradation in NLP benchmark performance caused by RLHF.

\paragraph{Key experimental results.}
(1) \textbf{The 1.3B InstructGPT outperformed the 175B GPT-3 in human evaluations}---
one of the most compelling experimental results in the post-training literature,
directly demonstrating that alignment can achieve qualitative leaps without increasing model scale.
(2) PPO-ptx maintained RLHF alignment effectiveness while
keeping NLP benchmark performance degradation to a minimal range,
whereas pure PPO caused significant performance decline.
(3) Pairwise agreement among annotators was approximately 73\%---
a figure that both validates that human preference judgments are consistent enough to support RLHF,
and suggests that reward model performance is upper-bounded by the human agreement ceiling.

\begin{solutionbox}
\textbf{The legacy of InstructGPT}:
Established the ``SFT $\rightarrow$ RM $\rightarrow$ PPO'' three-stage pipeline as the default paradigm for LLM post-training.
Subsequent models including Llama 2 \cite{touvron2023llama2}, Claude, and GPT-4 all adopted variants of this basic framework.
It also exposed the core bottleneck of RLHF:
the pipeline involves the coordinated training of four separate models
(SFT model, reward model, policy model, reference model),
resulting in extremely high engineering complexity and computational overhead.
\end{solutionbox}

\paragraph{Summary of the classic RLHF framework.}
From the prototype of Christiano et al.\ to the industrial-scale implementation of InstructGPT,
the RLHF framework underwent three key transitions:

\begin{enumerate}[leftmargin=2em]
  \item \textbf{From trajectory segments to full responses}:
    In NLP settings, reward is no longer summed over each timestep
    but is instead assigned as a single scalar score to the complete $(x, y)$ pair.
  \item \textbf{From real-time annotation to offline datasets}:
    Christiano et al.\ required human annotators to provide feedback online,
    whereas InstructGPT's RM was trained on pre-collected offline preference data.
  \item \textbf{From single-objective optimization to multi-objective balancing}:
    InstructGPT introduced KL penalty + pretraining loss mixing,
    establishing a multi-objective trade-off framework among helpfulness, safety, and capability retention.
\end{enumerate}

% ============================================================
\subsection{Limitations and Branches: Four Divergent Research Paths}
\label{sec:t2-branches}
% ============================================================

Although the classic RLHF framework proved effective, it exposed multiple core bottlenecks in practice.
These bottlenecks gave rise to four relatively independent research branches,
each attempting to address specific RLHF limitations from a different angle.

% --- 2.3.1 ---
\subsubsection{Branch 1: Reward Hacking and Overoptimization}
\label{sec:t2-reward-hacking}

\begin{problembox}
When the policy over-optimizes the reward model's output,
the RM score continues to rise while actual human preference evaluations begin to decline.
This is a direct manifestation of Goodhart's Law in RLHF:
``When a measure becomes a target, it ceases to be a good measure.''
\end{problembox}

\paragraph{Gao et al.'s scaling laws.}
\landmark{gao2022overoptimization} systematically studied the reward overoptimization phenomenon,
proposing scaling laws that describe the relationship between proxy RM scores and gold RM scores.

Let $d_{\text{KL}} = \text{KL}[\pi_\theta \| \pi_{\text{ref}}]$ measure the degree to which the policy deviates from the reference.
Gao et al.\ define $d = \sqrt{d_{\text{KL}}}$ as the normalized KL distance.
For Best-of-$N$ (BoN) sampling and RL (PPO) optimization, the gold reward (evaluated by a larger, more accurate RM) exhibits different functional forms with respect to $d$:

\begin{align}
\label{eq:pref:overopt-bon}
R_{\text{BoN}}^{\text{gold}}(d) &= d\left(\alpha_{\text{BoN}} - \beta_{\text{BoN}} \cdot d\right) \\[4pt]
\label{eq:pref:overopt-rl}
R_{\text{RL}}^{\text{gold}}(d) &= d\left(\alpha_{\text{RL}} - \beta_{\text{RL}} \cdot \log d\right)
\end{align}

\noindent
where the $\alpha$ term captures the signal correctly learned by the reward model (the initial reward growth rate),
and the $\beta$ term captures the rate at which reward model errors are exploited (the overoptimization term).
The key findings are as follows:

\begin{itemize}[leftmargin=2em]
  \item Both $\alpha$ and $\beta$ exhibit a \textbf{log-linear scaling} relationship with RM parameter count:
    larger RMs have higher $\alpha$ (capturing more true signal) and lower $\beta$ (harder to exploit),
    but overoptimization never completely vanishes.
  \item BoN is more \textbf{KL-efficient} than RL:
    under the same KL budget, BoN achieves higher gold reward.
    This implies that RL more readily exploits reward model vulnerabilities.
  \item KL penalty as a defense against overoptimization is necessary but insufficient:
    it can limit the growth of $d$, but cannot change the magnitude of $\beta$---
    that is, the intrinsic vulnerabilities of the reward model do not disappear due to KL constraints.
\end{itemize}

Gao et al.\ also systematically analyzed the four forms of Goodhart's Law as manifested in RLHF:
\textit{regressional} (regression error between proxy and gold reward),
\textit{extremal} (inconsistent ranking between proxy and gold in the distributional tails),
\textit{causal} (optimizing the proxy disrupts its causal relationship with the gold),
and \textit{adversarial} (the agent actively exploits vulnerabilities in the proxy).

\paragraph{Reward ensemble methods.}
To address overoptimization,
Coste et al.\ \cite{coste2023ensemble} proposed using reward model ensembles,
reducing single-RM bias through voting or averaging across multiple RMs.
However, subsequent work by Eisenstein et al.\ \cite{eisenstein2023ensemble} found that
while ensembles can mitigate overoptimization,
they cannot eliminate it entirely---when policy optimization is sufficiently strong,
shared vulnerabilities across all RMs in the ensemble can still be exploited.

\paragraph{Process Reward Model (PRM).}
The idea of process supervision did not originate with PRM.
Cobbe et al.\ \cite{cobbe2021gsm8k} first proposed training \textbf{outcome-based verifiers}
to judge whether the final answer of mathematical reasoning is correct, demonstrating on GSM8K that
verifier-guided search (Best-of-$N$) significantly outperforms single-sample generation.
Uesato et al.\ \cite{uesato2022process} then conducted the first systematic comparison of
\textbf{process-based feedback} (step-by-step annotation) versus \textbf{outcome-based feedback} (annotating only the final answer),
finding that process supervision achieves superior final-answer accuracy compared to outcome supervision,
and exhibits stronger discriminative power against ``false positive'' solutions (those with incorrect reasoning but coincidentally correct final answers).

Building on this foundation, \landmark{lightman2023prm} proposed a more systematic process reward model,
advancing reward granularity from outcome-level to step-level:
rather than assigning a single scalar reward to the entire response,
\textbf{each step} of the reasoning process is independently evaluated.

Specifically, in mathematical reasoning tasks,
a solution is decomposed into multiple steps $s_1, s_2, \ldots, s_n$,
and the PRM independently judges the correctness of each step:
\begin{equation}
\label{eq:pref:prm}
r_{\text{PRM}}(x, s_1, \ldots, s_n) = \prod_{i=1}^{n} P(s_i \text{ is correct} \mid x, s_1, \ldots, s_{i-1})
\end{equation}

Compared with the Outcome Reward Model (ORM)---which provides reward based solely on final answer correctness---
PRM offers finer-grained signals and can more precisely locate erroneous steps.
On the MATH dataset, PRM-guided Best-of-$N$ sampling achieved \textbf{78.2\%} accuracy ($N=1860$),
compared to 72.4\% for ORM, a significant improvement.

\begin{solutionbox}
\textbf{Key contributions of PRM}:
(1) Released the PRM800K dataset---containing approximately 800K step-level human labels,
covering 75K solutions and 12K mathematics problems;
(2) demonstrated that active learning achieves a \textbf{2.6$\times$} data efficiency improvement in PRM training;
(3) advanced reward modeling granularity from outcome-level to process-level,
an idea that subsequently had profound influence on reasoning verification (\crossref{4}{sec:t4}) and
process supervision in GRPO (\S\ref{sec:t2-grpo}).
\end{solutionbox}

% --- 2.3.2 ---
\subsubsection{Branch 2: RL-Free Methods---From RLHF to Direct Alignment}
\label{sec:t2-rlfree}

\begin{problembox}
The RLHF pipeline involves the coordinated training of four separate models (SFT model, reward model, policy, reference).
PPO training on LLMs is highly unstable and sensitive to hyperparameters,
and errors introduced by the reward model are amplified through RL optimization.
Can we bypass the explicit reward model and RL algorithm entirely,
and directly optimize the language model from preference data?
\end{problembox}

\paragraph{DPO: Direct Preference Optimization.}
\landmark{rafailov2023dpo} demonstrated through a key mathematical derivation
that the RL stage of RLHF can be completely eliminated---
the reward model can be \textbf{implicitly parameterized} by the language model itself.

The starting point of the \textit{derivation} is the standard RL objective of RLHF:
\begin{equation}
\label{eq:pref:rl-obj}
\max_{\pi_\theta}\; \mathbb{E}_{x \sim \mathcal{D},\; y \sim \pi_\theta(\cdot|x)} \left[
  r(x, y)
\right] - \beta\, \text{KL}\!\left[\pi_\theta(y|x) \| \pi_{\text{ref}}(y|x)\right]
\end{equation}

The core insight of DPO is that this KL-constrained optimization problem admits a \textbf{closed-form solution}.
For any reward function $r(x, y)$, the optimal policy is:
\begin{equation}
\label{eq:pref:optimal-policy}
\pi_r^*(y|x) = \frac{1}{Z(x)}\, \pi_{\text{ref}}(y|x) \exp\!\left(\frac{r(x, y)}{\beta}\right),
\quad Z(x) = \sum_y \pi_{\text{ref}}(y|x) \exp\!\left(\frac{r(x, y)}{\beta}\right)
\end{equation}

Taking the logarithm of \cref{eq:pref:optimal-policy} and rearranging, the reward can be expressed in terms of the policy and reference policy:
\begin{equation}
\label{eq:pref:reward-reparameterization}
r(x, y) = \beta \log \frac{\pi_r^*(y|x)}{\pi_{\text{ref}}(y|x)} + \beta \log Z(x)
\end{equation}

Substituting \cref{eq:pref:reward-reparameterization} into the Bradley-Terry preference model,
noting that the partition function $Z(x)$ cancels between the preferred and dispreferred responses,
yields the final DPO loss function:
\begin{equation}
\label{eq:pref:dpo-loss}
\boxed{
\mathcal{L}_{\text{DPO}}(\theta) = -\mathbb{E}_{(x, y_w, y_l) \sim \mathcal{D}} \left[
  \log \sigma\!\left(
    \beta \log \frac{\pi_\theta(y_w|x)}{\pi_{\text{ref}}(y_w|x)}
    - \beta \log \frac{\pi_\theta(y_l|x)}{\pi_{\text{ref}}(y_l|x)}
  \right)
\right]
}
\end{equation}

\paragraph{Gradient analysis of DPO.}
The gradient of the DPO loss carries deep intuitive meaning. Differentiating \cref{eq:pref:dpo-loss}:
\begin{equation}
\label{eq:pref:dpo-gradient}
\nabla_\theta \mathcal{L}_{\text{DPO}} = -\beta\, \mathbb{E}\!\left[
  \underbrace{\sigma\!\left(\hat{r}_\theta(x, y_l) - \hat{r}_\theta(x, y_w)\right)}_{\text{per-example weight}}
  \left(
    \nabla_\theta \log \pi_\theta(y_w|x) - \nabla_\theta \log \pi_\theta(y_l|x)
  \right)
\right]
\end{equation}
where $\hat{r}_\theta(x, y) \triangleq \beta \log \frac{\pi_\theta(y|x)}{\pi_{\text{ref}}(y|x)}$
is the \textbf{implicit reward}.

The weight term $\sigma(\hat{r}_\theta(x, y_l) - \hat{r}_\theta(x, y_w))$ in the gradient exhibits an adaptive property:
when the implicit reward already ranks $y_w$ above $y_l$ (i.e., $\hat{r}_\theta(x, y_w) \gg \hat{r}_\theta(x, y_l)$),
the sigmoid value approaches 0, reducing that sample's contribution to the gradient---
the model has already ``learned'' this preference pair and needs no further optimization.
Conversely, when the model still incorrectly prefers $y_l$, the sigmoid value approaches 1 and the gradient weight is maximal.
This \textbf{dynamic per-example importance weighting} is one of the keys to DPO's effectiveness.

\paragraph{Experimental validation.}
On the Anthropic HH dataset and TL;DR summarization task,
DPO achieved performance comparable to or better than PPO-based RLHF at lower computational cost.
More importantly, DPO's implementation requires only standard supervised learning infrastructure---
no value network, no online sampling, no PPO's complex clipping mechanism---
lowering the barrier to preference optimization from specialized RL research teams to any practitioner with SFT experience.

\paragraph{The DPO family tree.}
The simplicity and effectiveness of DPO spawned a large number of variants, forming a rich algorithmic family tree.
These variants primarily improve along the following dimensions:

\begin{itemize}[leftmargin=2em]
  \item \textbf{Preference model assumptions}:
    IPO \cite{azar2023ipo} relaxes the Bradley-Terry assumption,
    using a general representation of pairwise preferences
    $\Psi(r(x, y_w) - r(x, y_l)) \geq \Psi(0)$,
    and derives an MSE-form loss,
    avoiding DPO's overfitting problem under deterministic preferences.
    f-DPO \cite{wang2023fdpo} generalizes DPO to arbitrary $f$-divergence constraints.

  \item \textbf{Removing the reference model}:
    ORPO \cite{hong2024orpo} integrates the preference loss directly into the SFT objective,
    penalizing the generation probability of dispreferred responses via an odds ratio,
    eliminating the need for $\pi_{\text{ref}}$.
    SimPO \cite{meng2024simpo} uses average log-probability as the implicit reward:
    $\hat{r}_{\text{SimPO}}(x, y) = \frac{\beta}{|y|} \log \pi_\theta(y|x)$,
    likewise requiring no reference model.

  \item \textbf{Relaxed data requirements}:
    KTO \cite{ethayarajh2024kto}, based on Kahneman-Tversky prospect theory,
    requires only pointwise ``thumbs up / thumbs down'' labels
    rather than pairwise comparisons, substantially reducing data collection costs.

  \item \textbf{Theoretical unification}:
    GPO \cite{tang2024gpo} unifies DPO, IPO, SLiC-HF \cite{zhao2023slic}, and others
    within a framework parameterized by a general loss function.
    CPO \cite{xu2024cpo} and Smaug (DPO-Positive) \cite{pal2024smaug}
    mitigate DPO's tendency to reduce preferred response probability through additional SFT loss terms.
\end{itemize}

\evolution{RLHF: 4 models, RL training}{DPO: 2 models, supervised loss}{The reward model can be implicitly parameterized}

% --- 2.3.3 ---
\subsubsection{Branch 3: Online vs.\ Offline Preference Optimization}
\label{sec:t2-online-offline}

\begin{problembox}
DPO and its variants train on static, pre-collected preference data (offline),
but a distribution mismatch exists between the preference data distribution and the current policy's generation distribution.
As the policy is updated, the original preference data becomes increasingly off-policy,
causing the optimization signal to gradually lose effectiveness.
\end{problembox}

\paragraph{The necessity of on-policy data.}
The theoretical root of this problem lies in the classic on-policy vs.\ off-policy distinction in RL.
Tajwar et al.\ \cite{tajwar2024onpolicy} rigorously demonstrated that
even \textbf{suboptimal} on-policy data is superior to high-quality off-policy data.
Their core argument concerns distribution coverage---
off-policy data may fail to cover the response space that the current policy tends to generate,
which is precisely the region where policy improvement needs signals the most.

Song et al.\ \cite{song2024coverage} analyzed this problem from a more formal perspective,
proving that the sample complexity of online methods is superior to that of offline methods,
provided that online sampling offers sufficient coverage of the current policy distribution.

\paragraph{Iterative and online methods.}
Xiong et al.\ \cite{xiong2023iterative} proposed \textbf{Iterative DPO}:
in each round, the current policy generates new responses,
preferences are re-labeled via a reward model or AI judge,
and DPO training is performed on the new data.
This iterative approach approximates the effect of on-policy training
while retaining DPO's simple implementation.

Guo et al.\ \cite{guo2024oaif} proposed \textbf{OAIF} (Online AI Feedback),
using LLM-as-a-Judge to replace human annotators and an independent reward model,
achieving a fully online preference optimization pipeline---
at each step, the policy generates a response pair,
another LLM immediately judges preferences, and gradient updates are performed.

\paragraph{Self-play methods.}
Munos et al.\ \cite{munos2024nlhf} proposed
\textbf{Nash Learning from Human Feedback (NLHF)} from a game-theoretic perspective,
modeling preference optimization as a two-player constant-sum game
and seeking a Nash equilibrium policy.
Wu et al.\ \cite{wu2024sppo}'s \textbf{SPPO} (Self-Play Preference Optimization)
continued this line of thinking, using a self-play mechanism where the policy competes against its own historical versions,
progressively improving generation quality.

% --- 2.3.4 ---
\subsubsection{Branch 4: AI Feedback and Alternatives to Human Annotation}
\label{sec:t2-ai-feedback}

\begin{problembox}
The fundamental bottleneck of RLHF lies not in the algorithm but in the data:
high-quality human preference annotations are expensive, slow, and difficult to scale across all task types.
When model capabilities exceed annotator skill levels, human annotations may even introduce erroneous preference signals.
Can AI systems replace (or assist) human annotators?
\end{problembox}

\paragraph{Constitutional AI (CAI).}
\landmark{bai2022constitutional} proposed a systematic method
for AI systems to self-improve based on an explicit set of principles (a constitution).
The process consists of two stages:
\begin{enumerate}[leftmargin=2em]
  \item \textbf{Critique-Revision}:
    The model first generates a potentially harmful response, then is asked to critique its own response
    according to a principle from the constitution and generate a revised version.
    Multiple rounds of critique-revision iteration produce progressively improved responses.
  \item \textbf{RLAIF} (RL from AI Feedback):
    An AI model (rather than humans) provides preference judgments on response pairs,
    trains a reward model, and then performs RL optimization.
\end{enumerate}

\noindent
The core contribution of CAI is transforming alignment objectives from vague preferences implicit in human annotations
into an \textbf{explicit set of principles} that can be inspected and iterated upon.
This not only improves the interpretability and controllability of the alignment process,
but also provides a pathway toward scalable oversight
(\crossref{3}{sec:t3}).

\paragraph{Systematic comparison of RLAIF vs.\ RLHF.}
Lee et al.\ \cite{lee2023rlaif} systematically compared
the effectiveness of RLAIF and RLHF on summarization and helpful dialogue tasks,
finding that RLAIF achieved performance comparable to or even better than RLHF on most tasks.
This result suggests that when the AI annotator is sufficiently capable,
the marginal value of human annotation may be lower than expected.

\paragraph{Self-Rewarding Language Models.}
Yuan et al.\ \cite{yuan2024selfrewarding} pushed this idea to its extreme:
the language model simultaneously serves as both policy and reward model---
it generates responses, then uses its own LLM-as-a-Judge capability to evaluate response quality,
and performs iterative training based on self-assigned scores.
Experiments showed that multiple rounds of self-rewarding training can simultaneously improve
both the model's generation ability and its judgment ability,
forming a positive feedback loop of ``capability bootstrapping.''
However, this approach is upper-bounded by the model's own judgment capability---
when the model's judgments exhibit systematic biases, self-rewarding may amplify rather than correct these biases.

% ============================================================
\subsection{Convergence and Frontiers: New Balance Between Efficiency and Effectiveness}
\label{sec:t2-convergence}
% ============================================================

The four branching research paths described above have recently begun to exhibit convergence trends:
researchers seek methods
that possess both the on-policy characteristics and exploration capabilities of RL-based approaches,
while maintaining the simplicity and stability of DPO-like methods.
GRPO and RLOO are representative of this convergence trend.

% --- 2.4.1 ---
\subsubsection{GRPO: Group Relative Policy Optimization}
\label{sec:t2-grpo}

\landmark{shao2024grpo} proposed GRPO in DeepSeekMath,
with a core idea of elegant simplicity:
\textbf{no separate reward model or value network is needed---
only relative comparisons within a group of outputs from the same prompt.}

\paragraph{Algorithm details.}
For each prompt $x$, $G$ outputs $\{y_1, y_2, \ldots, y_G\}$ are sampled from the current policy $\pi_{\theta_{\text{old}}}$,
and each output receives a reward $r_i$ (which can come from an ORM, PRM, or rule-based verifiable reward).
The key step is \textbf{group-level reward normalization}:
\begin{equation}
\label{eq:pref:grpo-advantage}
\tilde{r}_i = \frac{r_i - \text{mean}(\{r_1, \ldots, r_G\})}{\text{std}(\{r_1, \ldots, r_G\})}
\end{equation}

A clipped objective similar to PPO is then used for optimization,
but with the advantage replaced by the group-normalized reward:
\begin{equation}
\label{eq:pref:grpo-obj}
\mathcal{L}_{\text{GRPO}}(\theta) = -\mathbb{E}_{x}\; \frac{1}{G} \sum_{i=1}^{G} \frac{1}{|y_i|} \sum_{t=1}^{|y_i|}
\left[
  \min\!\left(
    \rho_{i,t}\, \tilde{r}_i,\;
    \text{clip}(\rho_{i,t}, 1{-}\epsilon, 1{+}\epsilon)\, \tilde{r}_i
  \right)
  - \beta\, \text{KL}_t
\right]
\end{equation}
where $\rho_{i,t} = \frac{\pi_\theta(y_{i,t} | x, y_{i,<t})}{\pi_{\theta_{\text{old}}}(y_{i,t} | x, y_{i,<t})}$
is the token-level importance ratio,
and $\text{KL}_t$ is the token-level KL divergence penalty:
\begin{equation}
\text{KL}_t = \frac{\pi_{\text{ref}}(y_{i,t} | x, y_{i,<t})}{\pi_\theta(y_{i,t} | x, y_{i,<t})} - \log \frac{\pi_{\text{ref}}(y_{i,t} | x, y_{i,<t})}{\pi_\theta(y_{i,t} | x, y_{i,<t})} - 1
\end{equation}

\paragraph{A unified perspective on GRPO.}
Shao et al.\ proposed a highly illuminating \textbf{unified paradigm},
viewing SFT, RFT, DPO, PPO, and GRPO as special cases of the same objective function.
Specifically:
\begin{itemize}[leftmargin=2em]
  \item \textbf{SFT}: $G=1$, reward is always 1 (no selection).
  \item \textbf{RFT} (Rejection Sampling Fine-Tuning):
    $G>1$, only the sample with the highest reward is retained, equivalent to BoN + SFT.
  \item \textbf{DPO}: $G=2$, using the comparison signal from implicit rewards.
  \item \textbf{PPO}: $G=1$, advantage is estimated by a separate value network.
  \item \textbf{GRPO}: $G>1$, advantage is computed via group-relative normalization instead of a value network.
\end{itemize}

\noindent
This unified perspective reveals a profound connection:
\textbf{the essential difference among all preference optimization methods lies solely in
how the advantage function is estimated and how training data is sampled.}

\paragraph{Experimental results.}
In DeepSeekMath, GRPO combined with outcome reward + process reward
achieved \textbf{88.2\%} on GSM8K and \textbf{51.7\%} on MATH,
significantly outperforming SFT and RFT baselines.
GRPO later became one of the core training algorithms of DeepSeek-R1 \cite{deepseek2025r1}
(\crossref{4}{sec:t4}),
demonstrating powerful scaling capabilities on large-scale reasoning tasks.

\begin{solutionbox}
\textbf{Practical value of GRPO}:
Eliminates the additional memory overhead and training instability introduced by value networks,
bringing the resource requirements of RL-based preference optimization close to DPO levels,
while retaining the benefits of on-policy training and exploration.
\end{solutionbox}

% --- 2.4.2 ---
\subsubsection{RLOO and Other PPO Alternatives}
\label{sec:t2-rloo}

\paragraph{RLOO: REINFORCE Leave-One-Out.}
Ahmadian et al.\ \cite{ahmadian2024rloo} revisited the classic REINFORCE algorithm,
proposing a minimalist yet effective baseline estimation strategy:
for each prompt, $K$ responses $\{y_1, \ldots, y_K\}$ are sampled,
and the baseline for the $i$-th response is the average reward of the remaining $K-1$ responses:
\begin{equation}
\label{eq:pref:rloo}
b_i = \frac{1}{K-1} \sum_{j \neq i} r(x, y_j)
\end{equation}

This leave-one-out baseline provides a low-variance advantage estimate
without the need to train an additional value network.
Experiments showed that RLOO achieves performance comparable to PPO on RLHF tasks,
but with significantly lower implementation complexity and memory overhead.

\paragraph{Other PPO alternatives.}
Liu et al.\ \cite{liu2023rso}'s RSO (Rejection Sampling Optimization)
samples high-quality responses from the policy via rejection sampling,
then trains with a DPO-like loss,
occupying a conceptual middle ground between BoN and DPO.
Dong et al.\ \cite{dong2023raft}'s RAFT (Reward Ranked Fine-Tuning)
more directly combines rejection sampling with SFT:
a large number of responses are sampled, the top-$k\%$ by reward are retained,
and standard SFT training is performed on them.

\paragraph{RTO: DPO Meets PPO.}
Zhong et al.\ \cite{zhong2024rto} proposed a method that combines the strengths of both DPO and PPO:
using the implicit reward obtained from DPO training as a replacement for an independent reward model,
then performing PPO-style on-policy optimization,
seeking to balance DPO's simplicity with PPO's on-policy characteristics.

% --- 2.4.5 ---
\subsubsection{RLVR: Reinforcement Learning with Verifiable Rewards}
\label{sec:t2-rlvr}

In 2025, an important paradigm branch emerged in the preference optimization field:
\textbf{RLVR} (Reinforcement Learning with Verifiable Rewards),
which completely bypasses learned reward models
and directly uses rule-based automated verifiers as reward signals for RL training.
Although the idea of RLVR was already foreshadowed in R1-Zero \cite{deepseek2025r1}
(\crossref{4}{sec:t4-r1}),
DAPO \cite{yu2025dapo} systematized it into a reproducible open-source methodology.

\paragraph{The essence of the paradigm shift.}
The core assumption of traditional RLHF is that
human preferences need to be ``translated'' into optimizable signals through a reward model.
RLVR challenges this assumption---
\textbf{for tasks with automated verifiers (mathematics: answer correctness; programming: test pass rates),
rule-based rewards are more reliable, cheaper, and free from reward hacking compared to learned reward models}.
This insight deeply binds the methodology of Thread~2 to the reasoning tasks of Thread~4.

\paragraph{Systematic improvements in DAPO.}
\landmark{yu2025dapo} introduced four key improvements to GRPO:
\begin{enumerate}[leftmargin=2em]
  \item \textbf{Clip-higher}: A larger clip upper bound is used for positive advantage samples
    ($1 + \varepsilon_{\text{high}}$ instead of the standard $1 + \varepsilon$),
    encouraging the policy to more aggressively explore high-reward solution paths
    and addressing the excessive suppression of exploration caused by standard clipping;
  \item \textbf{Dynamic sampling}:
    Samples with identical rewards within a group (samples carrying no contrastive information) are dynamically filtered out,
    ensuring that every training batch contains effective contrastive signals;
  \item \textbf{Token-level loss}:
    The sequence-level policy gradient loss is decomposed into a token-level loss,
    providing finer-grained gradient signals and mitigating credit assignment issues in long sequences;
  \item \textbf{Removal of KL constraint}:
    The KL divergence penalty against the reference model is completely removed,
    allowing the policy to freely diverge from the initial distribution.
    For reasoning tasks, this ``unconstrained'' strategy yields significant gains,
    since the base model's initial reasoning distribution is typically far from optimal.
\end{enumerate}
On 7B models, DAPO improved AIME performance from $\sim$30\% with GRPO to $\sim$50\%.

\paragraph{Subsequent developments in RLVR.}
After DAPO, the RLVR paradigm rapidly spawned multiple variants:
\begin{itemize}[leftmargin=2em]
  \item \textbf{GSPO} (Group Sequence Policy Optimization) \cite{zheng2025gspo}:
    Further optimized the group-level sequence-wise policy update mechanism on top of DAPO,
    improving training stability on long-sequence reasoning tasks;
  \item \textbf{REINFORCE++} \cite{hu2025reinforcepp}:
    Proposed a streamlined and efficient REINFORCE variant that,
    through improved baseline estimation and gradient normalization strategies,
    achieves performance comparable to GRPO while maintaining a minimal implementation.
\end{itemize}

\paragraph{Theoretical perspective: Rewards as Labels.}
Zhai et al.\ \cite{zhai2026real} revisited RLVR from the perspective of classification learning,
pointing out that when rewards are binary and verifiable,
the RLVR training process can be understood as a special form of \textbf{classification with noisy labels}.
This theoretical perspective provides new analytical tools for studying the generalization properties of RLVR.

\begin{openproblembox}
\textbf{Boundaries of RLVR applicability}:
RLVR has achieved tremendous success on verifiable tasks such as mathematics and programming,
but these tasks inherently possess ``perfect reward functions.''
For open-ended tasks (dialogue, writing, reasoning-based QA),
can the RLVR paradigm be generalized?
One possible direction is to use LLM-as-judge as an ``approximate verifier,''
but this brings us back to the old problem of learned rewards.
Does RLVR's success imply that post-training methodology is splitting into two independent paradigms---
RLVR for verifiable tasks and RLHF for open-ended tasks?
\end{openproblembox}

% --- 2.4.6 ---
\subsubsection{2026: Theoretical Deepening and Methodological Breakthroughs}
\label{sec:t2-2026}

In early 2026, the field of preference optimization achieved important advances along three dimensions:
\textbf{repairs and theoretical unification of the DPO family},
\textbf{theoretical resolution of the online vs.\ offline debate},
and \textbf{paradigmatic innovation in reward modeling}.

\paragraph{Theoretical unification and key repairs of DPO.}
Raheja \& Pochhi \cite{raheja2026unification} provided a
\textbf{grand theoretical unification} of preference learning methods:
the seemingly diverse variants---DPO, IPO, KTO, SimPO, etc.---
can fundamentally be decomposed along three orthogonal axes:
(I) preference model assumptions, (II) regularization mechanisms, and (III) data distribution (online vs.\ offline).
This taxonomic framework provides practitioners with a systematic guide for method selection.

Regarding specific improvements to DPO,
HyPO \cite{hypo2026} identified a key failure mode of DPO---
\textbf{premature satisfaction}:
when the reference model favors the rejected response,
DPO prematurely weakens the gradient signal.
HyPO addresses this through \textbf{conditional reference regularization} (a one-line code modification),
achieving a 41.2\% average relative improvement on AlpacaEval 2.0.
PEPO \cite{barla2026pepo}, from a theoretical perspective,
provided the first proof that
\textbf{DPO's over-optimization can be provably avoided without knowing the data distribution}:
through a policy ensemble trained on disjoint data subsets,
it reduces sample complexity from dependence on all-policy concentrability
to dependence on single-policy concentrability alone.

\paragraph{Theoretical resolution of the on-policy question.}
\landmark{kim2026coverage} Kim et al.\ provided the
\textbf{first rigorous theoretical proof} that online preference learning is superior to offline.
The core conclusion is the \textbf{coverage improvement principle}:
with sufficient batch size,
each on-policy update moves the policy to a region of strictly better coverage,
thereby achieving exponential convergence.
More importantly, this work established a
\textbf{sharp sample complexity separation} between online and offline preference learning---
not only explaining why Iterative DPO and OAIF
consistently outperform vanilla DPO in practice,
but theoretically establishing the \emph{irreplaceability} of on-policy training.

\paragraph{Evolution of PPO alternatives.}
OPO \cite{wang2026opo} revealed a structural problem shared by both PPO and DPO:
\textbf{both confound sampling geometry and optimization geometry}.
By performing optimization in the Hilbert function space $L^2(\pi_k)$,
OPO decouples these two geometric structures,
resolving the numerical instability and gradient saturation problems in KL-regularized objectives.
R2VPO \cite{luo2026r2vpo} discovered that
PPO/GRPO's hard clipping \textbf{truncates high-reward but high-divergence ``eureka moments''}---
precisely those rare samples with the greatest learning value are cut off by gradient clipping.
R2VPO replaces hard clipping with a \emph{variance constraint} on the policy ratio,
achieving a 17\% relative gain with 50\% fewer rollouts.
He et al.\ \cite{he2026unifiedrlhf}
further revealed the implicit trade-off in RLHF between KL regularization (preventing reward hacking)
and ratio clipping (training stability),
proposing a unified regularization method---
whose final form reduces to a \textbf{weighted SFT loss},
yet consistently outperforms both RLHF and online preference learning.
This finding reignited the debate over ``is RL truly necessary?''

LLMdoctor \cite{llmdoctor2026} extended preference optimization
from training-time to \textbf{inference-time alignment}:
a small ``doctor'' model guides a frozen large ``patient'' model
through token-level flow-guided preference optimization,
achieving multi-dimensional preference alignment without retraining,
with effectiveness surpassing full DPO fine-tuning.

\paragraph{Paradigmatic innovation in reward modeling.}
Two works in 2026 may fundamentally change the data requirements for reward modeling.
\landmark{fan2026scalingrm} Fan et al.\ demonstrated that
reward models can be trained entirely from \textbf{document prefixes/suffixes in web corpora},
with \textbf{zero human annotation}.
Using only 11M math web data tokens,
they achieved up to 7.7 points of improvement on RewardBench v2,
with gains transferable across model families and scales.
If reward models can be trained without any preference annotations,
RLHF's data bottleneck would be fundamentally broken.

\landmark{young2026rlaif} Young proposed the \textbf{Latent Value Hypothesis}
to explain a long-standing puzzle: why does RLAIF work?
According to the data processing inequality, AI-generated preference signals should not be superior to human annotations,
yet RLAIF performs impressively in practice.
The hypothesis posits that the pretraining process encodes human values as \textbf{directions in representation space},
and constitutional prompts \emph{elicit} these latent values as preference judgments.
This theory unifies several scattered empirical findings:
refusal directions, low-rank safety subspaces, and RLAIF scaling.

Mix-GRM \cite{mixgrm2026} advanced the methodology of generative reward models:
decomposing the reasoning process into \textbf{Breadth-CoT} (multi-dimensional principle coverage)
and \textbf{Depth-CoT} (substantive judgment depth),
jointly optimizing both dimensions through SFT + RLVR,
and surpassing leading open-source RMs by 8.2\% across five benchmarks.
ActiveUltraFeedback \cite{melikidze2026activeuf}
introduced \textbf{uncertainty-based active learning} on the preference data collection side,
achieving or surpassing static baselines with only 1/6 of the annotation data.

\evolution{Preference data requires human annotation (RLHF) or AI annotation (RLAIF)}{Reward models can be trained unsupervised from web data (Fan et al.)}{Annotation cost is RLHF's core bottleneck}

% ============================================================
\subsection{Open Problems}
\label{sec:t2-open}
% ============================================================

Despite the rapid progress in preference optimization over the past two years,
the following key questions remain without satisfactory answers.

% --- 2.5.1 ---
\subsubsection{The Essential Difference Between DPO and RLHF}
\label{sec:t2-dpo-vs-rlhf}

DPO and RLHF should theoretically converge to the same optimal policy
(since the DPO loss is a transformation of the RLHF RL objective),
yet in practice their performance exhibits systematic differences.

Xu et al.\ \cite{xu2024dposuperior} systematically compared DPO and PPO,
finding that PPO generally outperforms DPO on large-scale models and complex tasks,
especially in scenarios requiring creative generation and open-ended tasks.
Ivison et al.\ \cite{ivison2024unpacking} analyzed the key factors
influencing their relative performance at a finer granularity,
including data quality, model scale, and training dynamics.

Rafailov et al.\ \cite{rafailov2024rtoq} reinterpreted DPO from the perspective of Q-functions,
pointing out that DPO's implicit reward $\hat{r}_\theta(x, y)$ can in fact
be understood as the value of a token-level Q-function at specific token positions.
However, DPO uniformly distributes the entire sequence's preference signal across all tokens,
whereas RLHF, through PPO's temporal credit assignment, can allocate rewards more precisely.
This may be the theoretical root cause of the performance gap between the two.

\begin{openproblembox}
\textbf{Core open problem}: Can the performance gap between DPO and RLHF be closed through better
offline preference optimization algorithms?
Or do on-policy training and explicit reward modeling possess irreplaceable advantages?
Recent empirical studies \cite{performance2025gap} suggest
that this gap may be jointly determined by credit assignment granularity and the degree of distribution shift.
\end{openproblembox}

% --- 2.5.2 ---
\subsubsection{Likelihood Displacement}
\label{sec:t2-displacement}

Razin et al.\ \cite{razin2024displacement} discovered a noteworthy phenomenon in DPO training:
\textbf{likelihood displacement}---during DPO training,
not only does the probability of the dispreferred response decrease,
but \textbf{the probability of the preferred response may also decrease},
albeit by a smaller magnitude.

This runs counter to intuitive expectations---
we expect DPO to increase the probability of preferred responses,
but in reality DPO's primary learning mechanism is
reducing the probability of dispreferred responses (creating a ``likelihood gap'')
rather than boosting the absolute probability of preferred responses.
When the probability of preferred responses also declines significantly (excessive displacement),
the model may lose useful generation capabilities.

\begin{openproblembox}
\textbf{Core open problem}: How can we design preference optimization algorithms
that widen the likelihood gap between preferred and dispreferred responses
while maintaining the absolute probability of preferred responses?
Methods such as CPO \cite{xu2024cpo} and Smaug \cite{pal2024smaug}
partially mitigate this issue by adding SFT loss terms,
but more principled solutions are still being explored.
\end{openproblembox}

% --- 2.5.3 ---
\subsubsection{Reward Model Generalization and Reliability}
\label{sec:t2-rm-generalization}

The generalization capability of reward models---whether they can still provide reliable preference signals
on prompt-response pairs outside the training distribution---is a fundamental open problem.

Zheng et al.\ \cite{zheng2023secrets1} and Wang et al.\ \cite{wang2024secrets2}
systematically studied key factors in reward model training,
including model architecture, data quality, and training objectives.
They found that reward model performance is highly sensitive to these factors,
and significant distribution shift problems exist---
an RM trained on one set of prompt types may severely degrade on novel types.
This problem is particularly acute in multimodal settings,
where vision-language preference signal annotation is more expensive and more subjective
(\crossref{5}{sec:t5-branch-a}).

A deeper question is: when the policy model's capabilities exceed those of the reward model,
does the reward model remain an effective training signal?
This is closely related to the problem of scalable oversight
(\crossref{3}{sec:t3})---
how to construct an evaluation system that can reliably assess model outputs that exceed its own capability level.

\begin{openproblembox}
\textbf{Core open problem}:
Are reward model generalization failure modes predictable?
Does PRM's process-level supervision fundamentally alleviate ORM's generalization problems,
or does it merely shift the problem to the reliability of step-level judgments?
How can we construct a reward signal that remains effective even when facing a superhuman policy?
\end{openproblembox}

\subsection{Chapter Summary}

\begin{figure}[!ht]
\centering
\begin{tikzpicture}[
  node distance=0.6cm and 1.2cm,
  every node/.style={font=\small},
  milestone/.style={rectangle, rounded corners, draw=thread2, fill=thread2!10,
    text width=4.2cm, minimum height=0.8cm, align=center, font=\small\bfseries},
  branch/.style={rectangle, rounded corners, draw=gray!60, fill=gray!5,
    text width=3.2cm, minimum height=0.7cm, align=center, font=\small},
  arrow/.style={-{Stealth[length=2.5mm]}, thick, thread2!70},
  brancharrow/.style={-{Stealth[length=2mm]}, gray!60, thick},
]
% Main timeline
\node[milestone] (christiano) {Christiano et al.\ (2017)\\Comparison-based RL};
\node[milestone, below=1.2cm of christiano] (instructgpt) {InstructGPT (2022)\\SFT$\rightarrow$RM$\rightarrow$PPO Pipeline};
\node[milestone, below=1.2cm of instructgpt] (dpo) {DPO (2023)\\Implicit Reward Optimization};

% Branches
\node[branch, below left=1.2cm and 0.5cm of dpo] (brA) {Branch A: Signal Enhancement\\Constitutional AI, PRM};
\node[branch, below right=1.2cm and 0.5cm of dpo] (brB) {Branch B: Offline Variants\\IPO, KTO, SimPO};

% Convergence
\node[milestone, below=2.8cm of dpo] (grpo) {GRPO (2024)\\Group Relative Policy Opt.};
\node[milestone, below=1.2cm of grpo] (dapo) {DAPO / RLVR (2025)\\Verifiable Reward Paradigm};
\node[branch, below=1.2cm of dapo] (conv) {Paradigm Divergence: Verifiable$\rightarrow$RLVR\\Open tasks$\rightarrow$RLHF/DPO};

% Arrows
\draw[arrow] (christiano) -- (instructgpt);
\draw[arrow] (instructgpt) -- (dpo);
\draw[brancharrow] (dpo) -- (brA);
\draw[brancharrow] (dpo) -- (brB);
\draw[arrow] (dpo) -- (grpo);
\draw[arrow] (grpo) -- (dapo);
\draw[brancharrow] (brA) -- (grpo);
\draw[brancharrow] (brB) -- (grpo);
\draw[arrow] (dapo) -- (conv);
\end{tikzpicture}
\caption{Thread 2 evolutionary path: from comparison-based reward learning to the RLVR paradigm divergence.
Blue nodes represent landmark works; gray nodes represent important branches.}
\label{fig:pref-evolution}
\end{figure}

\begin{table}[!ht]
\centering
\caption{Comparison of core methods in Thread 2. ``Needs RM'' indicates whether an independent learned reward model is required.}
\label{tab:pref-comparison}
\small
\begin{tabularx}{\textwidth}{l c c >{\raggedright\arraybackslash}X >{\raggedright\arraybackslash}X}
\toprule
\textbf{Method} & \textbf{Needs RM?} & \textbf{On/Offline} & \textbf{Core Mechanism} & \textbf{Limitation} \\
\midrule
RLHF/PPO & Yes & Online & PPO with KL penalty against $\pi_{\text{ref}}$; explicit reward model & 4 models; training instability \\
DPO & No & Offline & Implicit reward via $\log \frac{\pi_\theta}{\pi_{\text{ref}}}$; supervised loss & Distribution shift; likelihood displacement \\
IPO & No & Offline & MSE loss; relaxes Bradley-Terry assumption & Same offline issues as DPO \\
KTO & No & Offline & Pointwise labels only (thumbs up/down) & Weaker signal than pairwise \\
SimPO & No & Offline & Avg log-prob as reward; no $\pi_{\text{ref}}$ needed & No reference anchor \\
GRPO & Optional & Online & Group-normalized advantage; no value network & Multiple samples per prompt \\
DAPO & No (rule) & Online & RLVR: clip-higher + dynamic sampling & Verifiable tasks only \\
\bottomrule
\end{tabularx}
\end{table}

The evolutionary path of preference optimization can be summarized as follows:
starting from the comparison-based reward learning of Christiano et al.,
through the establishment of the standard paradigm via InstructGPT's industrial-scale three-stage pipeline,
and then continuously reducing method complexity and resource requirements through works such as DPO and GRPO,
while maintaining or improving alignment effectiveness.
The core tension along this evolutionary path lies in
\textbf{the trade-off between the richness of alignment signals and the simplicity of the training pipeline}---
RLHF provides the richest signals (on-policy + explicit reward + temporal credit assignment),
but at the cost of maximal engineering complexity;
DPO simplifies the pipeline to its extreme, but sacrifices on-policy characteristics and credit assignment precision;
GRPO and RLOO attempt to find new Pareto optima between these two extremes.
The emergence of the RLVR paradigm in 2025 further reshaped this landscape---
for verifiable tasks, methods like DAPO demonstrated that completely bypassing learned reward models
is not only feasible but superior,
leading to an emerging paradigm divergence in preference optimization methodology:
``RLVR for verifiable tasks, RLHF for open-ended tasks.''

\paragraph{Preview of the next chapter.}
Preference optimization makes model outputs better aligned with human preferences,
but the tension between ``helpfulness'' and ``safety'' remains a fundamental challenge---
a model that excessively pursues helpfulness may compromise on the safety dimension.
Thread~3 will delve into the evolution of safety alignment,
tracing the development path from HH-RLHF's exposure of this tension,
through Constitutional AI, the jailbreaking arms race, to representation engineering.
